{\color{red}
\FloatBarrier
\subsection*{What the three approximations say}

\paragraph{The reflector saving.} Water and iron behave as reflectors should. Both cores lose 14--20\% of their radius behind water and 18--34\% behind iron by theory (c), which because the mass goes as $R^3$ is a factor of more than three in critical mass: Pu-239 falls from 9.188 kg bare to 2.829 kg behind 10 mfp of iron. The two reflectors saturate differently, and the reason is $c_R$. Water absorbs ($c_R = 0.90$, $\nu_0 = 1.90$), so a neutron more than about two mean free paths deep never returns and the water columns are flat past $d = 3$. Iron barely absorbs ($c_R = 0.998$, $\nu_0 = 12.9$), so its return keeps improving out to $d = 10$ and beyond - iron starts as the worse reflector at $d=1$ and ends the better one at $d = 10$ under (a) and (b). Under (c) it is already ahead at $d = 1$.

\paragraph{The ordering of the three theories.} At every entry $R_c^{(a)} > R_c^{(b)} > R_c^{(c)}$. The first gap is inherited from the bare problem: the asymptotic $D_0(c_C)$ and $k_0(c_C)$ already give a smaller sphere than $D_0 = 1/3$ does, by 11\% for Pu-239. The second is the jump itself. Every pair here has $\mu_0(c_C) < \mu_0(c_R)$, so $g < 1$ - $0.766$ for Pu-239 behind water, $0.873$ for U-235 behind iron - because the multiplying core has the shorter relaxation length and hence the more forward-peaked angular flux, and so carries more current per unit density. A $g$ below one multiplies down the leakage term of \eqref{eq:reflfixed}, which is why (c) always predicts the smallest core: 7--12\% below (b) for water, 10--21\% for iron and 22--37\% for sodium. Part of that second gap is the change of criticality equation rather than the jump itself. Appendix~A separates the two.

\paragraph{The sodium columns of (a) and (b) are diffusion theory failing.} Every sodium radius in those two columns is \emph{larger} than the corresponding bare sphere, which no reflector can do. The condition is being solved correctly - \eqref{eq:reflcrit} was checked against an independent two-region spherical finite-volume $k$-eigenvalue solve, 3000 cells with harmonic-mean face diffusion coefficients across the material jump, and the two agree to better than $0.4\%$ on exactly these cases. The fault is in the model. Splitting the interface leakage of the Pu-239 core at $d = 10$ mfp into its two parts, $J/\phi = D_R/R_c + D_R/L$, the far-face return being $1/L = \Sigma_{t,R}\,\nu_0^{-1}\coth\!\left(\left[d+z_0\right]/\nu_0\right)$:


{\color{red}
\section*{Appendix A: what the change of criticality equation is worth}

Question 1 solves (a) and (b) from \eqref{eq:reflcrit} and (c) from \eqref{eq:reflfixed}. The
two are not the same equation with a different $\rho_{2/1}$ in it, so part of the gap between
the (b) and (c) columns is the jump and part is the change of equation. This appendix separates
them.

\paragraph{A.1 The gap is one bracket.} Carry the general $\rho_{2/1}$ and $j_{2/1}$ through
\eqref{eq:reflcrit} and collect its two $1/r$ terms, and it lines up against
\eqref{eq:reflfixed} term by term:
\begin{equation}
    \cot(k_0a) \;=\; \underbrace{\frac{1}{k_0a}\left(1 - \frac{\rho_{2/1}}{j_{2/1}}\frac{D_R}{D_C}\right)}_{\text{in \eqref{eq:reflcrit}, not in \eqref{eq:reflfixed}}} \;-\; \frac{\rho_{2/1}}{j_{2/1}}\frac{D_R/\nu_0}{D_C k_0}\coth\!\left(\frac{d + z_0(c_R)}{\nu_0(c_R)}\right)
    \label{eq:reflgap}
\end{equation}
with $D_R/D_C$ the ratio of the two \emph{dimensional} diffusion coefficients,
$D_{0,R}\Sigma_{t,C}/D_{0,C}\Sigma_{t,R}$. The underbraced bracket is the whole difference: it
carries the two $1/r$ curvature terms of \eqref{eq:reflcrit}, and \eqref{eq:reflfixed} does not
have it, because imposing the jump on $u = r\phi$ absorbs them. It is negative in every pair
here - $-0.12$ for U-235 behind water, where the two diffusion coefficients are within $35\%$
of each other, and $-3.28$ for Pu-239 behind sodium, where $D_R/D_C = 5.3$. Being negative it
lifts $a$, so dropping it is what pulls (c) below (b) by more than $\rho_{2/1} < 1$ alone
accounts for. Both equations still run from $+\infty$ at $a \to 0$ to $-\infty$ at
$a = \pi/k_0$, so the bracketing argument of Question 1 holds for either.

A consequence worth stating rather than discovering: (c) does not reduce to (b) as
$\mu_0(c_C) \to \mu_0(c_R)$. The two coincide only in the plane limit, where the bracket closes
by itself.

\paragraph{A.2 What it is worth.} Solving \emph{both} equations under \emph{both} sets of
parameters, for three entries of the Pu-239 table:

\begin{table}[htbp]
\centering
\begin{tabular}{ccccccc}
\toprule
 & & \multicolumn{2}{c}{(b) continuous asymptotic} & \multicolumn{2}{c}{(c) discontinuous asymptotic} & \\
reflector & $d$ [mfp] & \eqref{eq:reflcrit} & \eqref{eq:reflfixed} & \eqref{eq:reflcrit} & \eqref{eq:reflfixed} & reported \\
\midrule
Water  & 3  & 4.7499 & 4.4156 & 4.3321 & 4.1906 & 4.7499 / 4.1906 \\
Iron   & 10 & 4.4188 & 3.5484 & 4.1070 & 3.5037 & 4.4188 / 3.5037 \\
Sodium & 10 & 5.5286 & 3.5066 & 5.3080 & 3.4701 & 5.5286 / 3.4701 \\
\bottomrule
\end{tabular}
\caption{Critical radius in cm, Pu-239 core, under each equation and each parameter set. The
last column is the pair actually reported in Question 1. Bare asymptotic sphere: $5.1890$ cm.
Reading across the sodium row, the jump alone is worth $4\%$ and the change of equation a further $35\%$ - and
only \eqref{eq:reflfixed} keeps sodium below the bare sphere at all, which is the reason the
curvature-fixed form is the one used for (c).}
\end{table}

\paragraph{A.3 The price.} Matching $-D\,d[r\phi]/dr$ does not conserve $J$ across the
interface: it conserves $rJ + D\phi$, leaving a residual
$\left(\rho_{2/1}D_R - D_C\right)\phi_C/R_C^2$ of the same order as the geometric term it
discards. That is the content of the fix rather than an oversight - Fick's law is not the
true transport current near a curved interface either - and it is why the finite-volume
cross-check of the code, which is conservative by construction, reproduces \eqref{eq:reflcrit}
and can say nothing about \eqref{eq:reflfixed}.
}

\section*{Appendix: the coupled $P_N$ ODE system solved analytically}

The box scheme of Question 2 carries two errors, one in $N$ and one in $\Delta x$. Since
\eqref{eq:pnmatrix} is a linear system of coupled ODEs with constant coefficients, it can be
integrated in closed form at $k=1$ instead, with criticality imposed directly on the boundary
conditions. That removes the $\Delta x$ error altogether and leaves only the truncation in $N$,
which is what makes it an independent check on the mesh solution rather than a repetition of
it.

Split the $N+1$ moments into $M=\tfrac{N+1}{2}$ even and $M$ odd,
$\boldsymbol{\Phi}_{\text{even}}=(\phi_0,\phi_2,\dots)$ and
$\boldsymbol{\Phi}_{\text{odd}}=(\phi_1,\phi_3,\dots)$. Every derivative in \eqref{eq:pnrec}
changes the parity of its index, so the even-$n$ and odd-$n$ equations separate:
\begin{equation}
    \mathbf{A}\frac{d\boldsymbol{\Phi}_{\text{odd}}}{dx} + \mathbf{\Sigma}_0(k)\,\boldsymbol{\Phi}_{\text{even}} = 0,
    \qquad
    \mathbf{B}\frac{d\boldsymbol{\Phi}_{\text{even}}}{dx} + \boldsymbol{\Phi}_{\text{odd}} = 0,
    \qquad
    \mathbf{\Sigma}_0(k)=\text{diag}\!\left(1-\tfrac{c}{k},1,\dots,1\right)
\end{equation}
$\mathbf{A}$ and $\mathbf{B}$ being the two off-diagonal blocks of $\mathbf{\tilde A}$ under
that reordering, $\mathbf{\tilde A}=\left(\begin{smallmatrix}0&\mathbf{A}\\\mathbf{B}&0\end{smallmatrix}\right)$,
with rows $A_{m,\cdot}$ from $n=2m$ and $B_{m,\cdot}$ from $n=2m+1$. Substituting
$\boldsymbol{\Phi}_{\text{odd}}=-\mathbf{B}\,\boldsymbol{\Phi}_{\text{even}}'$ into the first
equation eliminates the odd moments and leaves a coupled Helmholtz system,
\begin{equation}
    \frac{d^2\boldsymbol{\Phi}_{\text{even}}}{dx^2} + \mathbf{K}^2(k)\,\boldsymbol{\Phi}_{\text{even}} = 0,
    \qquad
    \mathbf{K}^2(k) = -\,(\mathbf{A}\mathbf{B})^{-1}\mathbf{\Sigma}_0(k)
    \label{eq:helm}
\end{equation}
The minus sign is not cosmetic and is the easiest thing to lose: the elimination produces
$\mathbf{A}\mathbf{B}\,\boldsymbol{\Phi}_{\text{even}}''=\mathbf{\Sigma}_0\boldsymbol{\Phi}_{\text{even}}$,
and moving that to the form \eqref{eq:helm} flips it. At $N=1$ it reads $\phi_0''+3(c-1)\phi_0=0$,
the classic diffusion buckling $B=\sqrt{3(c-1)}$ of the table in Question 1 - with the sign dropped
one gets $\cosh$ where $\cos$ belongs and no critical thickness at all.

\paragraph{The spectrum of $\mathbf{K}^2$.} Since
$\mathbf{\tilde A}^2=\left(\begin{smallmatrix}\mathbf{AB}&0\\0&\mathbf{BA}\end{smallmatrix}\right)$ and the
eigenvalues of $\mathbf{\tilde A}$ are the $N+1$ roots $\pm\mu_i$ of $P_{N+1}$, the matrix
$\mathbf{A}\mathbf{B}$ has the $M$ eigenvalues $\mu_i^2>0$ and is similar to a positive definite
matrix. By Sylvester's law of inertia the eigenvalues of \eqref{eq:helm} then carry the signs of
$-\mathbf{\Sigma}_0$, which for $c>1$ at $k=1$ has one negative and $M-1$ positive entries:
\textbf{exactly one eigenvalue $\lambda_1^2=B_0^2>0$, and $M-1$ eigenvalues
$\lambda_j^2=-\kappa_j^2<0$}. The one oscillatory mode is the fundamental. The rest are the
boundary layers, and at $c=1$ the positive one collapses to $B_0=0$, an infinite critical slab.

\paragraph{Assembling the criticality condition.} Diagonalise
$\mathbf{K}^2(1)=\mathbf{V}\mathbf{\Lambda}\mathbf{V}^{-1}$. The solution even about $x=0$,
which satisfies \eqref{eq:pnsym} identically because
$\boldsymbol{\Phi}_{\text{odd}}=-\mathbf{B}\boldsymbol{\Phi}_{\text{even}}'$ vanishes wherever
$\boldsymbol{\Phi}_{\text{even}}'$ does, is
\begin{align}
    \boldsymbol{\Phi}_{\text{even}}(x) &= C_1\mathbf{v}_1\cos(B_0x)
        + \sum_{j=2}^{M} C_j\mathbf{v}_j\cosh(\kappa_jx) \\
    \boldsymbol{\Phi}_{\text{odd}}(x) &= C_1(\mathbf{B}\mathbf{v}_1)B_0\sin(B_0x)
        - \sum_{j=2}^{M} C_j(\mathbf{B}\mathbf{v}_j)\kappa_j\sinh(\kappa_jx)
\end{align}
$M$ free constants against the $M$ Marshak conditions \eqref{eq:marshak}. Splitting the
coefficient matrix $c_{mn}$ into its even and odd columns
$\mathbf{M}_{\text{even}},\mathbf{M}_{\text{odd}}$ and evaluating at $x=a/2$ gives $M$ homogeneous
equations $\mathbf{H}(a/2)\mathbf{C}=\mathbf{0}$ with
\begin{align*}
    \mathbf{H}_{\cdot1} &= \mathbf{M}_{\text{even}}\mathbf{v}_1\cos\!\left(\tfrac{B_0a}{2}\right)
        +\mathbf{M}_{\text{odd}}(\mathbf{B}\mathbf{v}_1)\,B_0\sin\!\left(\tfrac{B_0a}{2}\right) \\
    \mathbf{H}_{\cdot j} &= \mathbf{M}_{\text{even}}\mathbf{v}_j\cosh\!\left(\tfrac{\kappa_ja}{2}\right)
        -\mathbf{M}_{\text{odd}}(\mathbf{B}\mathbf{v}_j)\,\kappa_j\sinh\!\left(\tfrac{\kappa_ja}{2}\right),
        \qquad j\ge2
\end{align*}
and a non-trivial flux exists only where
\begin{equation}
    \boxed{\;\det\mathbf{H}(a/2)=0\;}
    \label{eq:hdet}
\end{equation}
$(a/2)_{\text{crit}}$ is its smallest positive root. Two numerical points: the $\cosh$ and $\sinh$
columns overflow once $\tfrac12\kappa_ja\gtrsim700$, so each is divided by its own $\cosh(\tfrac12\kappa_ja)$,
which leaves $\mathbf{v}_j$ and $-(\mathbf{B}\mathbf{v}_j)\kappa_j\tanh(\tfrac12\kappa_ja)$ and cannot
move a zero of the determinant, a column scaling being a non-zero factor of it. And the
determinant itself spans many orders of magnitude across the bracket, so it is the sign change
that is tracked, not the value.


\paragraph{Against the box scheme.} The two solutions are the only free check on each other,
since they share nothing but the matrices of \eqref{eq:pnmatrix} and
\eqref{eq:marshak}: one iterates a discretised operator, the other never discretises space at
all. They agree to between one and three parts in $10^{6}$, and that residue is the box scheme's
$\Delta x^2$ and nothing else. Refining the mesh at $c=1.5$, $N=5$ against the analytic
$0.612134509$:

\begin{center}
\begin{tabular}{cccccc}
\toprule
cells $J$ & 25 & 50 & 100 & 200 & 400 \\
\midrule
error [mfp] & $6.61\cdot10^{-5}$ & $1.65\cdot10^{-5}$ & $4.13\cdot10^{-6}$ & $1.03\cdot10^{-6}$ & $2.58\cdot10^{-7}$ \\
ratio & - & 4.00 & 4.00 & 4.00 & 4.00 \\
\bottomrule
\end{tabular}
\end{center}

Each halving of $\Delta x$ divides the error by $4.00$ to three digits, which is the second
order claimed in Question 2 measured rather than asserted, and it confirms that the boundary rows are
as accurate as the interior ones - a first-order boundary closure would have shown up here as
a ratio of $2$.
}
